\documentclass[11pt]{article}

\usepackage[margin=1in]{geometry}
\usepackage[T1]{fontenc}
\usepackage{amsmath,amssymb,amsthm,mathtools}
\usepackage{booktabs}
\usepackage{array}
\usepackage{xcolor}
\usepackage{listings}
\usepackage[protrusion=true,expansion=false]{microtype}
\usepackage{enumitem}
\usepackage[hidelinks,colorlinks=true,linkcolor=black,citecolor=blue!45!black,urlcolor=blue!45!black]{hyperref}

\theoremstyle{definition}
\newtheorem{definition}{Definition}[section]
\newtheorem{remark}{Remark}[section]
\theoremstyle{plain}
\newtheorem{proposition}{Proposition}[section]

\definecolor{rhokw}{rgb}{0.10,0.25,0.55}
\definecolor{rhocom}{rgb}{0.45,0.45,0.45}

\lstdefinestyle{rho}{
  basicstyle=\ttfamily\footnotesize,
  breaklines=true,
  columns=fullflexible,
  keepspaces=true,
  showstringspaces=false,
  commentstyle=\color{rhocom}\itshape,
  morecomment=[l]{//},
  keywordstyle=\color{rhokw}\bfseries,
  morekeywords={new,in,for,contract,match,if,else,true,false,Nil},
  frame=single,
  framesep=4pt,
  rulecolor=\color{black!30},
  xleftmargin=6pt,
  aboveskip=8pt,belowskip=8pt
}
\lstset{style=rho}

\newcommand{\rcalc}{rho calculus}
\newcommand{\sig}[1]{\ensuremath{s_{#1}}}
\newcommand{\wrap}[2]{\ensuremath{\{#1\}_{#2}}}
\newcommand{\redto}{\ensuremath{\rightsquigarrow}}

\title{\textbf{Financial Contract Combinators in the\\ Cost-Accounted Rho Calculus}\\[4pt]
\large An executable, resource-safe settlement language --- and a firmer footing for DeFi}

\author{L. G. Meredith\\ \small F1R3FLY.io \quad\textemdash\quad \texttt{lgreg.meredith@f1r3fly.io}}

\date{June 2026}

\begin{document}
\maketitle

\begin{abstract}
Peyton Jones, Eber and Seward showed that a small set of combinators can describe an
unbounded family of financial contracts, and that a compositional denotational semantics
can value them. Their language describes \emph{what a contract is} and \emph{what it is
worth}; it does not say how a contract \emph{executes} on a shared, adversarial, resource-metered
substrate, nor what guarantees that execution enjoys. We close that gap by realising the
combinator language in the cost-accounted \rcalc: the reflective higher-order process calculus
that is the formal core of Rholang, equipped with the cost endofunctor that folds resource
accounting into the grammar. Reflection supplies the substrate --- contracts, observables and
obligations are all processes, and names are quoted processes --- while cost accounting supplies
three things a settlement language needs and a denotational semantics cannot: \emph{located
authority} over funds, \emph{atomicity} of multi-step settlement, and a \emph{linear resource
proof} that a deployment is funded \emph{before} any step fires. We re-express five instrument
families in this calculus --- the core combinators, loans, a Dutch auction, cash-settled futures,
and a book of fourteen power-market instruments --- and in doing so identify a correction to the
original cost presentation (token stacks must be \emph{located}, because parallel composition is
associative--commutative) and a structural finding (the combinator \texttt{and} silently conflates
a simultaneous atomic swap with a bundle of independently-scheduled obligations; the acceptance gate
forces the two apart). We argue that the resulting discipline addresses, by construction, several
of the failure modes that have dogged decentralised finance --- stranded funds, ambient authority,
gas exhaustion mid-transaction, and post-hoc conflict resolution --- that it lightens the post-2008
regulatory burden by making the audit trail native and compliance a proof discharged before execution
rather than an audit performed after it, and that these same properties make it a credible internal
settlement substrate for incumbent finance. We close by looking forward
to the spatial-behavioral types that the calculus's own logic generator, OSLF, makes available as a
machine-checkable discipline for financial computing.
\end{abstract}

\tableofcontents

\section{Introduction}

A financial contract is an agreement to exchange cash flows contingent on time and on observable
quantities. Peyton Jones, Eber and Seward~\cite{SPJ2000} made the seminal observation that the
sprawling catalogue of traded instruments --- swaps, futures, caps, floors, options of every
flavour --- is generated by a handful of combinators, and that giving those combinators a
compositional denotational semantics lets one \emph{value} an arbitrary contract by structural
recursion. Their combinators (\texttt{zero}, \texttt{one}, \texttt{give}, \texttt{and}, \texttt{or},
\texttt{scale}, \texttt{truncate}, \texttt{get}, \texttt{anytime}) and the derived forms built from
them (\texttt{zcb}, \texttt{european}, \texttt{american}) remain the canonical answer to the question
``what is the right vocabulary for financial contracts?''

That work answers two questions --- what a contract \emph{is}, and what it is \emph{worth} --- and
deliberately leaves a third untouched: how a contract \emph{executes}. On a single trusted machine,
execution is unremarkable. On a distributed ledger --- where execution is concurrent, adversarial,
and metered, where a half-completed settlement can strand value, where the right to move funds is a
capability that must not leak, and where running out of fuel mid-transaction is a real and frequent
event --- execution is the entire problem. This is the regime of decentralised finance (DeFi), and
it is precisely the regime in which the combinator language has the most to offer and has been least
applied.

We supply the missing layer by realising the combinator language in the \emph{cost-accounted}
\rcalc~\cite{CAR2026}, a calculus obtained from the reflective higher-order \rcalc~\cite{MeredithRadestock2005a}
--- the formal core of Rholang~\cite{Rholang} --- by applying the cost endofunctor of~\cite{CGC2026}.
The thesis of this note is simple to state:

\begin{quote}
\emph{The combinator language, interpreted in the cost-accounted \rcalc, is not merely a description
language with a valuation semantics; it is an executable settlement language whose executions are
resource-safe and atomic by construction.}
\end{quote}

We make the case constructively, by re-expressing five instrument families and reading off what the
apparatus buys in each.\footnote{The full implementation --- rho source for all five instrument
families re-expressed in the cost-accounted calculus --- is open-sourced in the F1R3Fi
repository~\cite{F1R3Fi}.} Two findings emerge along the way that are of independent interest. First, a
correction: the original \rcalc{} presentation of cost accounting left token stacks floating in the
parallel-composition bag, and because that composition is associative--commutative this is an
\emph{ambient-authority} leak; the fix is to \emph{locate} every stack on an unforgeable channel ---
in this calculus a location is just a name --- so that authority is possession of the name, not
adjacency in the bag (\S\ref{sec:located}). Second, a structural observation: the combinator \texttt{and} is overloaded,
denoting both a simultaneous bilateral swap and a bundle of independently-scheduled obligations, and
the cost calculus's acceptance gate forces these apart because the gate's scope \emph{is} the unit of
financial atomicity (\S\ref{sec:patterns}).

\paragraph{Outline.} \S\ref{sec:background} recalls the \rcalc, the combinators, and the cost
construction. \S\ref{sec:realization} is the realisation: the contract protocol, located funding,
and the key features that distinguish a settlement language from a description language.
\S\ref{sec:instruments} surveys the five families and the four cost-accounting patterns they exhibit.
\S\ref{sec:defi} argues the case for DeFi. \S\ref{sec:types} looks ahead to spatial-behavioral types.
\S\ref{sec:conclusion} concludes.

\section{Background}\label{sec:background}

\subsection{The reflective higher-order \rcalc}

The \rcalc~\cite{MeredithRadestock2005a} is a process calculus in which \emph{names are quoted
processes}. Its grammar is
\[
P, Q \;::=\; 0 \;\mid\; \texttt{for}(y \leftarrow x)\,P \;\mid\; x!(Q) \;\mid\; P \mid Q \;\mid\; {*}x,
\qquad x, y \;::=\; @P,
\]
with input \texttt{for}$(y\leftarrow x)P$, output $x!(Q)$, parallel composition $P \mid Q$,
the quotation $@P$ of a process to a name, and the dereference ${*}x$ recovering the process a name
quotes. Communication substitutes the quoted datum for the bound name. Reflection --- the
$@\,/\,{*}$ pair --- is what makes the calculus higher-order without a separate function space, and it
is what lets contracts, observables, oracles and obligations all be \emph{processes}, exchanged as
\emph{names} on channels and run by dereference. This is the substrate on which Rholang and the
F1R3FLY platform are built.

\subsection{The combinators}

A contract, in the realisation below, is a channel \texttt{C} that, when sent an acquisition message,
discharges the rights and obligations it denotes. Following~\cite{SPJ2000} the primitives are
\texttt{zero} (no rights or obligations), \texttt{one}$(k)$ (acquire one unit of currency $k$),
\texttt{give}$(c)$ (reverse the parties of $c$), \texttt{and}$(c_1,c_2)$ (acquire both),
\texttt{or}$(c_1,c_2)$ (the holder chooses one), \texttt{scale}$(o,c)$ (scale $c$ by an observable
$o$), \texttt{truncate}$(t,c)$ (let $c$ expire at $t$), and \texttt{get}$(c)$ (acquire $c$ at its
horizon). An \emph{observable} is a time-varying quantity both parties agree on, modelled as a
persistent channel that returns its current value. From these one derives the zero-coupon bond
$\texttt{zcb}(t,x,k) = \texttt{scale}(x,\texttt{get}(\texttt{truncate}(t,\texttt{one}(k))))$, the
European option $\texttt{european}(t,u) = \texttt{get}(\texttt{truncate}(t, u \mathbin{\texttt{or}}
\texttt{zero}))$, and so on, exactly as in the source.

\subsection{Cost accounting as a construction on the calculus}

Rholang gates every deployment behind a \emph{phlogiston} balance tied to a signature. The cost
papers~\cite{CAR2026,CGC2026} show that this gating is not an ad hoc runtime extension but the image
of a generic endofunctor $C$ on a category of interactive process theories. The construction
adjoins, to the host calculus, three sorts and a discipline:

\begin{itemize}[leftmargin=1.4em,itemsep=2pt]
\item \emph{Signatures as names.} The name sort is extended, $x,y ::= @T \mid s$, so a signature $s$
is itself a channel. A signature both \emph{funds} the phlogiston of a communication and
\emph{authorises} the underlying value movement; there is no separate translation
layer~\cite[\S5.1]{CAR2026}. Signatures compose: $s_1 * s_2$ is joint (multi-signature) authority.
\item \emph{Token stacks as first-class terms.} A stack $S ::= () \mid s :: S$ is a temporal monoid;
its head is the next token to spend. A balance is a stack; minting and delegation are sends of
stacks on signature channels.
\item \emph{Wrapping by construction.} Every continuation is re-sorted to a wrapped-term sort: a
signed thunk $\wrap{P}{s}$ that fires only by spending a matching token. ``No redex fires without
being unwrapped'' is then a sorting invariant preserved by reduction~\cite[\S3]{CGC2026}, not a
property to be re-established by traversal; metering is lazy, charging for work performed rather than
exposed.
\end{itemize}

The single communication rule of the host is replaced by a gated family, of which the base case is
\[
\wrap{\texttt{for}(y\leftarrow x)\,T \mid x!(U)}{s} \;\mid\; s :: S
\;\;\redto\;\;
T\{@U/y\} \;\mid\; S,
\]
``a communication fires only when a matching token is present, and the consumed token is released as
residual.''~\cite[\S4.6]{CAR2026} Two consequences organise everything below. First, the
\texttt{for}-comprehension \emph{is} the transaction primitive: its binding specification is the
transaction, the substitution is the witness, and application to the continuation is the state
update; a single communication is atomic by construction~\cite[\S8.1]{CAR2026}. Second, \emph{financial}
atomicity --- that a multi-step operation completes entirely or not at all --- is \emph{not} a theorem
of the rewrite rules but is secured by a static analysis at deployment-acceptance time: a
\emph{linear resource proof} verifying $\Sigma_s \ge \Delta_s$, supply at least demand, for every
signature, \emph{before} any step executes~\cite[\S8.5]{CAR2026}.

\section{Realising the Combinators}\label{sec:realization}

\subsection{The cost-accounted contract protocol}

A contract becomes a channel acquired by
\[
\texttt{C}!(\,s_H,\ s_{Cp},\ \mathit{hAddr},\ \mathit{cAddr},\ \mathit{fund},\ \mathit{res}\,),
\]
where $s_H$ and $s_{Cp}$ are the holder and counterparty signatures (names that are simultaneously
capabilities and authorisations), $\mathit{hAddr}$, $\mathit{cAddr}$ are the parties' ledger
addresses, $\mathit{fund}$ is a \emph{located funding slot} (below), and $\mathit{res}$ returns
$(\texttt{true},\texttt{Nil})$ or $(\texttt{false},\mathit{err})$. The combinator \texttt{give}
swaps $(s_H,s_{Cp})$ and $(\mathit{hAddr},\mathit{cAddr})$ together: swapping \emph{which signature
funds} each leg is exactly the two-sided surface of the interaction cut~\cite[\S4.8.2]{CAR2026} made
operational. The payment primitive \texttt{one} is funded by the payer, drawing one token; the
structural combinators forward the funding slot to their sub-contracts. Thus \emph{cost incidence
follows economic incidence}: the party who pays the cash flow pays the phlogiston that moves it ---
a property that is simply inexpressible in a model with a single fee-paying deployer.

\subsection{Located token stacks: a correction}\label{sec:located}

The first realisation attempt exposes a defect in the original cost presentation. There, a token
stack $s :: S$ is written into the ambient process soup and consumed by whatever communication needs
it. But parallel composition is associative--commutative: the soup is an unordered bag, and
\emph{any} process in it is adjacent to \emph{every} stack in it. Taking adjacency as the right to
spend is \emph{ambient authority}~\cite{Miller2006} --- the antithesis of an object-capability
discipline, and a latent double-spend.

The fix is to \emph{locate} every stack, and in the rho calculus a location is nothing more exotic
than a \emph{name}. To locate a token stack is simply to hold it on a channel: the stack becomes a
message $c!(\,s :: S\,)$ on a channel $c$, and the right to spend from it is exactly the right to
receive on $c$. Because channels minted by \texttt{new} are unforgeable, possession of $c$ is a
genuine object capability --- a process that was never handed $c$ cannot name it, cannot receive on
it, and so cannot touch the stack, even though, the bag being associative--commutative, it sits in the
same soup. The name does the work that position cannot: it says \emph{who may draw} (possession of
$c$), while the signature at the head of the stack says \emph{which token is spent} (the guard).
Authority is possession of a name, never adjacency. In the realisation every fuel stack lives on such
a channel; none floats free.

This is the concrete, rho-calculus reading of the general \emph{located resource stack}
$S(I,-)$ of~\cite[\S10.4]{CGC2026}, and the abstraction collapses pleasantly here. In general the
location $I$ may be carried either structurally, by rigid position, or nominally, by an explicit name;
but an associative--commutative combinator like $\mid$ destroys position, so in this setting the
location can \emph{only} be a name. What the general construction calls an ``interaction surface'' is,
for us, just a channel. The same mechanism is the \emph{funding slot} of~\cite[\S5.7]{CAR2026}: an
open contract --- an auction whose bidders are unknown in advance --- creates a fresh channel and
accepts fuel on it, funded by whoever holds the channel, with no need to name them at compile time.

\subsection{The acceptance gate and financial atomicity}

The motivating hazard of the cost papers is a payment decomposed into a debit and a credit that funds
the debit but not the credit, stranding value in transit~\cite[\S8.2]{CAR2026}. Every multi-step
settlement in the combinator language has this shape, and the apparatus closes it the same way: a
multi-step acquisition first \emph{reserves} its entire demand from a located slot, all-or-nothing,
before any leg fires. Listing~\ref{lst:reserve} shows the reservation gate; it is the operational
content of the linear resource proof $\Sigma_s \ge \Delta_s$, decided before execution. If the
reservation fails the slot is left untouched and nothing moves; if it succeeds, every step draws from
the committed purse and cannot stall for lack of fuel.

\begin{lstlisting}[caption={The acceptance gate: an atomic reservation from a located slot, the operational form of the linear resource proof.},label={lst:reserve}]
// reserve(L, sig, n, ret): remove the first n tokens of `sig`
// from located slot L, all-or-nothing; either (true, purse)
// with a committed private purse, or L is left UNTOUCHED.
contract reserve(@L, @sig, @n, ret) = {
  for (@stack <- @L) {
    new takeCh, peel in {
      contract peel(@rem, @k, @acc, pret) = {
        match k {
          0 => pret!((true, acc, rem))
          _ => match rem {
            [] => pret!((false, Nil, Nil))
            [hd ...tl] => match hd == sig {
              true  => peel!(tl, k - 1, acc ++ [hd], *pret)
              false => pret!((false, Nil, Nil)) } } } } |
      peel!(stack, n, [], *takeCh) |
      for (@result <- takeCh) {
        match result {
          (true, taken, remainder) => {
            @L!(remainder) | new P in { P!(taken) | ret!((true, *P)) } }
          (false, _, _) => { @L!(stack) | ret!((false, "insufficient fuel")) }
        } } } }
}
\end{lstlisting}

\begin{remark}[Two atomicities]
Database-transaction atomicity --- a single communication either fires or does not --- is a theorem
of the rewrite rules. Financial-transaction atomicity --- a multi-step settlement completes or does
not begin --- is secured by the gate above, at the deployment boundary. The combinator tree is the
unit over which demand is summed; acceptance is the unit of financial atomicity. This is the precise
sense in which the description language becomes a settlement language.
\end{remark}

\subsection{The combinator \texttt{and} as an atomic join}

The combinator \texttt{and} acquires both sub-contracts. Naively encoded, its two legs are funded
independently, so one leg's transfer may commit while the other stalls --- a half-settled swap with
no rollback. The cost calculus offers the principled remedy: fund the join under the compound
authority $s_1 * s_2$ and reserve the joint demand atomically, so that, by the join cost
schema~\cite[\S4.8]{CAR2026}, the redex either gathers a full token's worth of every participant's
authority or does not fire. ``The partial-funding hazard \dots is structurally impossible within a
join.'' For a same-funder bundle (a fixed-payment strip) the joint reservation is from a single
located purse; for a cross-party swap (a loan sale, or the bilateral settlement below) it is a
two-slot reservation that funds both parties or neither.

\subsection{Data-dependent demand: conservative bound and refund}

The choice combinators (\texttt{or}, \texttt{european}) and the temporal combinators
(\texttt{get}, \texttt{truncate}) branch on data --- the holder's choice, or the block height ---
so which interactions fire, and hence how many tokens of which signature are drawn, depends on the
run. For these the demand is not a fixed sum. The discipline is conservative: reserve the worst case
over branches and refund the unforced remainder after the run~\cite[\S8.6]{CAR2026}. Over-funding is
a benign inefficiency reconciled by a refund; under-funding would reintroduce the stranding hazard,
so we never under-fund. The realisation tracks each combinator's demand as a conservative bound
($\max$ over branches for choice; the live branch for time) and the deployment boundary refunds what
the run did not force.

\subsection{Cost is not amount}

A subtlety that the prism metaphor of~\cite[\S5.6]{CAR2026} makes precise: phlogiston counts
\emph{funded rendezvous}, not the magnitude of value moved. Scaling a contract by a large observable
multiplies the currency transferred but leaves the fuel demand at one --- one transfer, one token,
whatever its size. The static funding proof concerns the count of rendezvous, not the REV moved;
conflating them (as a gas-per-byte intuition invites) is a category error the apparatus forbids.

\section{Five Families, Four Patterns}\label{sec:instruments}

We re-expressed five instrument families~\cite{F1R3Fi}: the core combinators; loans (short, coupon and amortising,
with an atomic loan sale); a Dutch auction; cash-settled futures; and a book of fourteen power-market
instruments (PPA, futures, swing, interruptible supply, capacity, REC forward, virtual PPA, demand
response, frequency regulation, spinning reserve, forward-curve strip, spark spread, weather
derivative, congestion revenue right). Across all five families the instruments fall into four
cost-accounting patterns, summarised in Table~\ref{tab:patterns}.

\begin{table}[t]
\centering\small
\begin{tabular}{@{}p{0.21\textwidth}p{0.34\textwidth}p{0.37\textwidth}@{}}
\toprule
\textbf{Pattern} & \textbf{Cost-accounting treatment} & \textbf{Instruments} \\
\midrule
(A) Bilateral cash settlement & Net the two opposing legs into one transfer; reserve one token from
each party, refund the unused side & Futures, virtual PPA, spark spread, CRR, weather derivative \\
\addlinespace
(B) Single-funder fixed leg & One funded transfer, demand $1$; \texttt{give} swaps the funder &
\texttt{one}, \texttt{zcb}, REC forward, capacity, PPA period \\
\addlinespace
(C) Conservative-bound option & \texttt{or}/\texttt{european}: demand $=\max$ over branches; writer
pre-funds the exercise case, refunded on lapse & Swing, frequency regulation, interruptible, option
legs of DR / spinning reserve \\
\addlinespace
(D) Time-separated bundle & Legs at different blocks and/or funded by different parties; each leg its
own funded deployment under concurrent acceptance & Strips (PPA, swing, freq, forward curve);
premium$+$option pairs (DR, spinning reserve); loan disburse$+$repay \\
\bottomrule
\end{tabular}
\caption{The four patterns that cover all the realised instruments.}
\label{tab:patterns}
\end{table}

\subsection{Bilateral settlement and the netting incentive}

The futures contract is the sharpest case. Its denotation is a bilateral obligation --- the
counterparty pays the holder the spot price, the holder pays the counterparty the agreed price ---
written $\texttt{and}(\texttt{scale}(\mathit{spot},\texttt{one}(1)), \texttt{give}(\texttt{one}(F)))$.
As two gross transfers in opposite directions, funded by different parties, this is the partial-funding
hazard at its worst. Cost accounting makes the gross form strictly worse on both axes it tracks: two
transfers can half-settle at the ledger layer, and they cost two tokens. The cost-accounted form
therefore \emph{nets}: sample the difference at expiry and perform a \emph{single} transfer in the
appropriate direction (Listing~\ref{lst:net}). A single transfer is atomic by construction and costs
one token regardless of magnitude. The settlement direction is data-dependent, so we reserve one token
from each party and refund the unused side. Netting is thus not merely an optimisation but the form
the apparatus rewards: it collapses a bilateral swap to a single atomic step and removes the
half-settlement hazard entirely. Five of the fourteen power instruments --- the genuine contracts for
differences --- settle through this one mechanism, parameterised by a (free) net observable.

\begin{lstlisting}[caption={Netted bilateral settlement: one transfer in the data-dependent direction; atomic by construction.},label={lst:net}]
// at expiry T, net = (what holder receives) - (what holder pays)
match net == 0 {
  true  => res!((true, Nil))                 // flat: nothing moves
  false => match net > 0 {
    true  => { spend!(cpPurse, *sp) | ...     // cp pays holder net
               transfer!(RevVault, cAddr, hAddr, net, sC, *res) }
    false => { spend!(holderPurse, *sp) | ... // holder pays cp -net
               transfer!(RevVault, hAddr, cAddr, 0 - net, sH, *res) }
  } }
\end{lstlisting}

\subsection{The overloading of \texttt{and}}\label{sec:patterns}

The power book reveals that \texttt{and} silently denotes two different things. In the futures it is a
\emph{simultaneous} swap: both legs settle at the same instant and must be jointly atomic. In a
demand-response contract it bundles a premium paid \emph{now} (grid to data centre) with an option
settled \emph{later} (data centre to grid, on exercise) --- two obligations at different times, funded
by different parties. The denotational semantics is indifferent to the distinction; the cost calculus
cannot be, because the acceptance gate's scope \emph{is} the unit of financial atomicity. Fusing a
time-separated bundle into a single atomic gate would over-reserve fuel and wrongly couple two
obligations that should stand or fall independently. The realisation therefore splits \texttt{and}
into an atomic join (Pattern A/C settlements) and a bundle returned as a list of independently
acquirable legs (Pattern D), each its own funded deployment under \emph{concurrent
acceptance}~\cite[\S2.3]{CAR2026}. That this distinction is invisible to valuation but forced by
execution is, we think, a small vindication of the thesis: the settlement semantics sees structure
the denotational semantics cannot.

\section{A Firmer Footing for DeFi}\label{sec:defi}

Decentralised finance has repeatedly stumbled over a small set of failure modes that are not failures
of cryptography or consensus but of \emph{execution discipline}~\cite{Atzei2017}. The cost-accounted
combinator language addresses each of them not by patching but by construction.
Table~\ref{tab:defi} summarises the correspondence.

\begin{table}[t]
\centering\small
\begin{tabular}{@{}p{0.30\textwidth}p{0.63\textwidth}@{}}
\toprule
\textbf{DeFi failure mode} & \textbf{Cost-accounted remedy} \\
\midrule
Stranded funds / half-settled multi-step transactions & The acceptance gate funds the whole
transaction or rejects it before any step fires; netting collapses bilateral swaps to a single
atomic transfer (\S\ref{sec:realization}). \\
\addlinespace
Ambient authority / unauthorised draws on pooled funds & Located resource stacks: fuel and value sit
on unforgeable surfaces; authority is possession of the surface, never adjacency in the
AC bag (\S\ref{sec:located}). \\
\addlinespace
Gas exhaustion mid-execution & A deployment is accepted only if a linear resource proof certifies it
is fully funded; an accepted deployment cannot run out of fuel. Over-estimates are refunded. \\
\addlinespace
Reentrancy & The unit of state update is a single \texttt{for}-comprehension, atomic by construction;
there is no partial intermediate state to re-enter, and capabilities are not ambient. \\
\addlinespace
Post-hoc conflict resolution / MEV from ordering & Concurrent acceptance over disjoint token pools:
deployments signed by different signatures cannot conflict and compose without merge analysis;
competing deployments are competing linear-proof obligations, at most one of which
succeeds~\cite[\S8.8]{CAR2026,Daian2020}. \\
\addlinespace
Opaque, unauditable contract logic & Contracts are declarative compositions of a fixed combinator set
with a denotational valuation~\cite{SPJ2000} and an operational settlement semantics; the same term
is valued, executed, and (\S\ref{sec:types}) typed. \\
\bottomrule
\end{tabular}
\caption{DeFi failure modes and their structural remedies in the cost-accounted realisation.}
\label{tab:defi}
\end{table}

The deeper point is methodological. Smart-contract platforms typically offer a Turing-complete
imperative language and a gas meter bolted on outside it, leaving atomicity, authority and resource
sufficiency as the programmer's responsibility --- which is to say, as the auditor's nightmare. The
cost-accounted \rcalc{} inverts this: resource accounting is \emph{inside} the grammar (``valuable
friction at the syntactic level''~\cite[\S5.3]{CAR2026}), authority is an object-capability discipline
carried by unforgeable names, and atomicity is either a theorem (single step) or a statically checked
proof obligation (multi-step). Building a financial DSL --- the combinator language --- on \emph{this}
substrate means the instruments inherit those guarantees rather than re-litigating them. DeFi does not
need a more clever contract; it needs a substrate on which the obvious contract is already safe.

\section{Regulatory Compliance and Internal Deployment}\label{sec:compliance}

The 2008 financial crisis was, in large part, a crisis of opacity: derivative exposures that no party
could see whole, contracts no system could value in aggregate, and obligations whose interlocking was
discovered only as they failed. The regulatory response built a standing apparatus of reporting,
recordkeeping, clearing, and capital discipline whose cost the industry now carries continuously. In
the United States the Dodd--Frank Act~\cite{DoddFrank2010} (Title~VII) mandated central clearing of
standardised swaps, real-time public reporting, and the reporting of every swap to a registered swap
data repository; in the European Union, EMIR~\cite{EMIR2012} imposed parallel trade-reporting and
clearing obligations together with margin requirements for uncleared derivatives, and MiFID~II /
MiFIR~\cite{MiFIDII2014} added exhaustive transaction-reporting and recordkeeping duties; Basel
III~\cite{BaselIII2011} tightened risk-based capital, leverage, and liquidity ratios; and the
BCBS--IOSCO uncleared-margin rules made collateralisation of non-cleared positions mandatory. The
domain these regimes police --- swaps, futures, options, structured payments --- is exactly the domain
of the combinator language. The cost-accounted substrate supports their obligations along two axes: a
\emph{record} that is native rather than reconstructed, and a \emph{proof} discharged before a contract
executes.

\paragraph{The record axis.} A distributed ledger is, by construction, a complete, timestamped,
tamper-evident history of every deployment and settlement. Here that record is unusually faithful,
because a contract is not opaque bytecode but a declarative combinator term with a denotational
valuation and an operational settlement trace: the record \emph{is} the contract, together with the
exact sequence of funded rendezvous that settled it. The audit trail that swap-data
reporting~\cite{DoddFrank2010}, EMIR trade reporting~\cite{EMIR2012}, and MiFID~II transaction
reporting~\cite{MiFIDII2014} demand at significant standing cost is, on this substrate, a by-product of
execution rather than a separate reporting pipeline bolted alongside it. Counterparty identity ---
the contribution of the post-crisis Legal Entity Identifier regime to transparency --- is carried
directly by the signatures-as-names that already fund and authorise each leg.

\paragraph{The proof axis.} The deeper opportunity is to discharge compliance \emph{statically},
before a contract is accepted, rather than auditing it afterwards. Because contracts are declarative
terms with formal semantics, and because the calculus already checks one resource property --- the
linear funding proof of \S\ref{sec:realization} --- at the deployment gate, the same gate can carry
regulatory proof obligations. Several map cleanly:
\begin{itemize}[leftmargin=1.4em,itemsep=2pt]
\item \emph{Collateralisation and margin.} The uncleared-margin requirement that a position be backed
by sufficient collateral is, almost verbatim, the located-purse resource-sufficiency proof of
\S\ref{sec:realization}: a contract that cannot prove it is funded does not execute. Margin compliance
is the funding proof under another name.
\item \emph{Authorised-party gates (KYC/AML).} A compound signature
$\{\,\mathit{transfer}\,\}_{s_{\mathrm{party}} * s_{\mathrm{compliance}}}$ makes a compliance officer's
co-signature a \emph{structural} precondition of settlement~\cite[\S2.4]{CAR2026}: the transfer cannot
fire without it. Authorisation is a capability the term carries, not a check performed after the fact.
\item \emph{Position and exposure limits.} The spatial types of \S\ref{sec:types} bound how
obligations may compose and how much of a resource a sub-contract may claim --- the type-level form of
a position limit or a concentration cap.
\end{itemize}
This shifts compliance from after-the-fact audit --- the expensive, error-prone posture the post-2008
regime largely institutionalised --- to a proof checked before any value moves: the regulatory analogue
of the resource-sufficiency guarantee. As with the type discipline generally (\S\ref{sec:types}), the
record axis is available today, while the full proof axis builds on the type layer still in progress;
the substrate is designed to carry both.

\paragraph{Internal deployment for incumbent finance.} None of these properties --- atomic settlement,
capability-confined authority, a native audit trail, pre-trade compliance proofs, concurrency without
post-hoc reconciliation --- depends on a public, permissionless deployment. They are equally available,
and arguably land first, inside a regulated institution running a private or permissioned shard as an
internal settlement-and-compliance engine: the same combinator contracts, the same guarantees, without
the regulatory uncertainty that still surrounds public decentralised finance. For a bank or a fintech
the value proposition is concrete and quantifiable --- the audit trail and pre-trade compliance proofs
attack a known, measured cost line --- which is often where new infrastructure earns its first
production deployment. The arrival of dedicated crypto-asset regimes such as the EU's
MiCA~\cite{MiCA2023} makes a substrate with compliance built in, rather than bolted on, especially
timely. The cost-accounted rho calculus is thus not only a firmer footing for open DeFi but a credible
internal infrastructure for the very institutions that the post-2008 reforms now bind.

\section{Related Work: Marlowe and the Substrate Question}\label{sec:related}

The closest prior art is Marlowe~\cite{LamelaThompson2018}, the financial-contract DSL developed for
Cardano in the same Peyton Jones--Eber lineage as the present work~\cite{SPJ2000}. Marlowe shares our
premise --- that financial contracts deserve a small, analysable, special-purpose language rather than
a general imperative one --- and pursued analysability aggressively: the language is deliberately
finite and non-Turing-complete, so that the behaviours of a contract can be enumerated and checked
exhaustively. It is instructive to ask why an approach so close in spirit did not become the substrate
for open decentralised finance, because the answer locates precisely where our bet differs.

The difficulty was not in the contract language but in the substrate beneath it. Cardano's Extended
UTXO model~\cite{EUTXO2020} places application state in unspent outputs, and an output may be consumed
by at most one transaction per block. Any \emph{contended, shared} state --- an order book, a
liquidity pool, an auction, any genuinely multi-party venue --- therefore becomes a serialisation
point: concurrent participants produce conflicting transactions, all but one fail, and applications
recover throughput only by interposing off-chain ``batchers'' that order and submit on users' behalf.
But most of open DeFi \emph{is} contended shared state, and batching reintroduces exactly the
centralisation and ordering discretion --- a privileged sequencer, a private mempool --- that a public
ledger was meant to dissolve~\cite{Daian2020}. The substrate fought the workload.

The rho calculus begins from the opposite commitment. It is a process calculus; parallel composition
is primitive, not an optimisation; RSpace is a concurrent tuplespace; and the cost construction's
governing motivation is to replace run-to-completion with \emph{concurrent acceptance} over disjoint
token pools~\cite[\S2]{CAR2026}: deployments signed by different signatures provably cannot conflict
and compose without post-hoc merge analysis, while deployments competing for the same funds are
competing linear-proof obligations of which at most one succeeds~\cite[\S8.8]{CAR2026}. Where EUTXO
makes shared-state concurrency the hard case to be worked around, the cost-accounted rho calculus
makes it the default case, decided structurally at acceptance time. This is the single strongest
reason to expect a different outcome, and it is independent of the contract language sitting on top.

The two designs also recover safety by opposite routes, and this is the second difference. Marlowe
buys exhaustive analysability by sacrificing expressiveness: a finite, non-Turing-complete language
cannot host the open-ended, composable, user-defined instruments that drove adoption on more
expressive platforms. We keep a universal substrate --- the rho calculus is Turing-complete --- and
recover safety \emph{structurally} rather than by restriction: atomicity by construction, located
object-capability authority, the linear resource proof checked before execution, and the
spatial-behavioral type discipline of \S\ref{sec:types}. The aim is Marlowe-grade guarantees without
Marlowe's expressiveness ceiling, with resource accounting inside the grammar rather than bolted on as
an external fee model. We are candid that this is in part a promissory note: the operational
guarantees are in hand in the realisation above, but the type-system layer that would make them
statically checkable is the work of \S\ref{sec:types}, still ahead of us. A process calculus also pays
for its concurrency with a larger reasoning burden than a finite language --- interleaving and
nondeterminism are real --- which is exactly why we regard the behavioural type discipline as
load-bearing rather than ornamental. Other efforts to give contracts a process-calculus footing, such
as BitML for Bitcoin~\cite{BartolettiZunino2018}, share the conviction that contracts deserve a
calculus; they differ in target and do not address the concurrent-settlement problem that is our focus.

Finally, candour about traction itself. Whether a contract language is adopted is not a purely
technical verdict; network effects, liquidity depth, and timing dominate, and Cardano's broader DeFi
ecosystem was thin and late relative to where capital had already pooled. We do not claim a technical
argument settles an economic outcome. The defensible claim is narrower and, we think, sound: the
cost-accounted rho calculus removes the specific \emph{architectural} reason a Peyton Jones--Eber
contract language could not host real decentralised finance on its predecessor substrate --- it makes
the obvious financial contract safe, concurrent, and composable at once, where EUTXO could deliver at
most two of the three.

\section{Looking Forward: Spatial-Behavioral Types}\label{sec:types}

The guarantees above are secured operationally --- by located slots, the acceptance gate, and the
conservative-bound discipline. The natural next step is to make them \emph{types}: machine-checkable
assertions a contract carries and a checker verifies. The calculus already generates the logic for
this. OSLF~\cite{MeredithRadestock2005b} auto-generates, from a graph-structured theory, a
Hennessy--Milner-style logic~\cite{HennessyMilner1985} that is adequate for bisimulation; run on the
\emph{cost-accounted} theory it yields a \emph{graded} logic whose modalities are indexed by the
signature consumed~\cite[\S7,\S11]{CGC2026}. Two connective families suffice to express a token
discipline as a type:
\begin{itemize}[leftmargin=1.4em,itemsep=2pt]
\item \emph{Spatial} types $K(\varphi_1,\varphi_2)$ read parallel composition as a separating
conjunction in the sense of separation and spatial
logics~\cite{OHearn2001,Reynolds2002,CairesCardelli2003}: a term satisfies $K(\varphi_1,\varphi_2)$
when it splits into disjoint parts satisfying $\varphi_1$ and $\varphi_2$. This is the bookkeeping of
ownership and duplication --- whether two obligations hold disjoint funds, whether a token is copied.
\item \emph{Modal} types $\langle K(\dots)\rangle\varphi$ are the graded modalities of consumption: a
term can take a funded interaction, after which $\varphi$ holds.
\end{itemize}
Together they express the lattice of resource disciplines --- linear (spend exactly once), affine,
relevant, copyable --- as \emph{checkable} assertions, decidable when stacks are finite and located.

To see the two families at work --- and that the $K$ they instantiate need not be the same one ---
take the netted future of \S\ref{sec:instruments} --- the bilateral contract for differences whose
\texttt{and} is a simultaneous swap (\S\ref{sec:patterns}). Its fuel sits on two located surfaces
(\S\ref{sec:located}): a holder purse
drawable by $\sig{H}$ and a counterparty purse drawable by $\sig{Cp}$. Instantiate the \emph{spatial}
$K$ as parallel composition, $\mathbin{\mid}$, and write $\mathsf{purse}_s$ for ``a located stack
keyed to $s$ sits here.'' The ownership type is a separation,
\[
  \Phi_{\mathrm{sep}} \;\equiv\; \mathsf{purse}_{\sig{H}} \mathbin{\mid} \mathsf{purse}_{\sig{Cp}},
\]
satisfied when the term splits into two disjoint purses sharing no sub-term: the two parties' funds
provably do not alias --- the static no-double-spend. Instantiate the \emph{behavioural} $K$,
by contrast, as the all-or-nothing \emph{join} cut of~\cite[\S4.8]{CAR2026} --- a \emph{different}
constructor, the settlement rendezvous rather than the spatial bag --- and grade its modality by the
signature it draws. The atomicity type is a single compound-graded diamond,
\[
  \Phi_{\mathrm{settle}} \;\equiv\; \langle\,\mathrm{join}\,\rangle_{\sig{H}*\sig{Cp}}\;\mathsf{Flat},
\]
read: the contract takes one funded step whose demand is the compound token $\sig{H}*\sig{Cp}$ --- one
cell from each party, drawn jointly --- after which the position is $\mathsf{Flat}$ (settled, nothing
stranded). Because that compound cell may not be partially consumed~\cite[\S4.8.5]{CAR2026}, the
diamond cannot half-fire. The two predicates read the two faces of the one GSLT
constructor~\cite[\S4]{CGC2026}: the spatial $K=\mathord{\mid}$ records \emph{what is near what} (the
purses are apart), the behavioural $K=\mathrm{join}$ records \emph{what is spent} (one compound step
settles them). The point is sharpened by the contract that fails to type: the \emph{gross}, un-netted
future admits no single such diamond, only the staged chain
$\langle\mathrm{cut}\rangle_{\sig{Cp}}\langle\mathrm{cut}\rangle_{\sig{H}}\,\mathsf{Flat}$, whose
reachable mid-state between the two diamonds is exactly the half-settled, stranded position. Netting
is what makes $\Phi_{\mathrm{settle}}$ available, and the half-settling contract does not type-check
against it.

For financial computing the payoff is concrete and, we believe, considerable. The future above is the
contract-for-differences case --- $\Phi_{\mathrm{settle}}$ types its bilateral atomicity,
$\Phi_{\mathrm{sep}}$ its no-double-spend by separation --- and the same two families reach further. An option
writer's obligation can be typed so that the worst-case exercise is provably funded for the life of
the option. A loan's repayment schedule can carry a per-instalment resource-sufficiency proof,
factored by located purse into independent conjuncts~\cite[\S12.2]{CGC2026}, so that funding a
hundred-period amortisation is a hundred small local checks rather than one global one. In each case
the property that the present realisation secures operationally becomes a type the term carries and a
checker discharges --- the financial analogue of memory safety. We regard this as the most promising
direction the present work opens: not safer contracts written more carefully, but a type discipline
under which the unsafe contract does not type-check.

This sits within a longer arc. The same spatial/temporal structure that underwrites the cost
calculus --- the spatial monoid $K$ of what is near what, the temporal monoid $::$ of what is spent
next --- is the structure on which the calculus's authors have noted suggestive correspondences to
physical field theories~\cite[\S13]{CGC2026}. We do not lean on those here; we note only that a
financial settlement layer, a behavioural type system, and a process-theoretic account of resource
flow are, in this setting, three readings of one object.

\section{Conclusion}\label{sec:conclusion}

We have realised the Peyton Jones--Eber--Seward combinator language in the cost-accounted \rcalc, and
in doing so turned a description-and-valuation language into an executable settlement language whose
executions are resource-safe and atomic by construction. Reflection supplies the substrate on which
contracts, observables and obligations are uniformly processes; the cost endofunctor supplies located
authority over funds, atomicity of multi-step settlement, and a linear resource proof checked before
execution. Re-expressing five instrument families --- combinators, loans, a Dutch auction, futures,
and a fourteen-instrument power book --- surfaced a correction (token stacks must be located, because
parallel composition is associative--commutative) and a structural finding (the acceptance gate forces
apart the simultaneous swap and the time-separated bundle that \texttt{and} conflates). We argued that
the discipline addresses, by construction, the execution-level failure modes that have dogged DeFi,
that its native audit trail and pre-execution compliance proofs lighten the post-2008 regulatory
burden and make the substrate a credible internal settlement engine for regulated institutions,
and we sketched the spatial-behavioral type system that the calculus's own logic generator makes
available as the next layer of robustness. A description language tells you what a contract is worth.
A settlement language on this substrate tells you, before a single token moves, that it will settle.

\section*{Acknowledgements}
This note develops joint work conducted in extended dialogue with Anthropic's Claude, which served as
a sounding board throughout --- re-expressing the combinator families in the cost-accounted calculus,
surfacing the located-stack correction and the overloading of \texttt{and}, and pressing on the
atomicity and resource-sufficiency arguments. The ideas and any errors are the author's.

\begin{thebibliography}{99}\small

\bibitem{SPJ2000}
S. Peyton Jones, J.-M. Eber, J. Seward.
\newblock Composing contracts: an adventure in financial engineering (functional pearl).
\newblock In \emph{ACM SIGPLAN International Conference on Functional Programming (ICFP)}, Montreal, 2000.

\bibitem{PJEber2003}
S. Peyton Jones, J.-M. Eber.
\newblock How to write a financial contract.
\newblock In J. Gibbons and O. de Moor, eds., \emph{The Fun of Programming}, Palgrave Macmillan, 2003.

\bibitem{MeredithRadestock2005a}
L. G. Meredith, M. Radestock.
\newblock A reflective higher-order calculus.
\newblock \emph{Electronic Notes in Theoretical Computer Science} 141(5):49--67, 2005.

\bibitem{MeredithRadestock2005b}
L. G. Meredith, M. Radestock.
\newblock Namespace logic: a logic for a reflective higher-order calculus.
\newblock In \emph{Trustworthy Global Computing (TGC)}, LNCS 3705:353--369, Springer, 2005.

\bibitem{CAR2026}
L. G. Meredith.
\newblock Cost-accounted rho calculus: a spectral decomposition of phlogiston.
\newblock F1R3FLY.io, 2026.

\bibitem{CGC2026}
L. G. Meredith.
\newblock Continued interactive GSLTs and the cost endofunctor: a construction one level up from
cost-accounted Rholang.
\newblock F1R3FLY.io, 2026.

\bibitem{Rholang}
L. G. Meredith et al.
\newblock Rholang specification.
\newblock F1R3FLY.io / RChain Cooperative, 2017--2026.

\bibitem{MeTTaIL}
L. G. Meredith.
\newblock MeTTaIL (Rholang 1.4): a language for defining languages.
\newblock F1R3FLY.io, 2025--2026.

\bibitem{Milner1999}
R. Milner.
\newblock \emph{Communicating and Mobile Systems: the $\pi$-Calculus}.
\newblock Cambridge University Press, 1999.

\bibitem{HennessyMilner1985}
M. Hennessy, R. Milner.
\newblock Algebraic laws for nondeterminism and concurrency.
\newblock \emph{Journal of the ACM} 32(1):137--161, 1985.

\bibitem{Girard1987}
J.-Y. Girard.
\newblock Linear logic.
\newblock \emph{Theoretical Computer Science} 50(1):1--101, 1987.

\bibitem{Reynolds2002}
J. C. Reynolds.
\newblock Separation logic: a logic for shared mutable data structures.
\newblock In \emph{IEEE Symposium on Logic in Computer Science (LICS)}, 55--74, 2002.

\bibitem{OHearn2001}
P. W. O'Hearn, J. C. Reynolds, H. Yang.
\newblock Local reasoning about programs that alter data structures.
\newblock In \emph{Computer Science Logic (CSL)}, LNCS 2142:1--19, Springer, 2001.

\bibitem{CairesCardelli2003}
L. Caires, L. Cardelli.
\newblock A spatial logic for concurrency (part I).
\newblock \emph{Information and Computation} 186(2):194--235, 2003.

\bibitem{Miller2006}
M. S. Miller.
\newblock \emph{Robust Composition: Towards a Unified Approach to Access Control and Concurrency
Control}.
\newblock PhD thesis, Johns Hopkins University, 2006.

\bibitem{StayMeredith2017}
M. Stay, L. G. Meredith.
\newblock Representing operational semantics with enriched Lawvere theories.
\newblock In \emph{Higher-Dimensional Rewriting and Applications (HDRA)}, 2017. arXiv:1704.03080.

\bibitem{deBruijn1972}
N. G. de Bruijn.
\newblock Lambda calculus notation with nameless dummies.
\newblock \emph{Indagationes Mathematicae} 34(5):381--392, 1972.

\bibitem{Atzei2017}
N. Atzei, M. Bartoletti, T. Cimoli.
\newblock A survey of attacks on Ethereum smart contracts (SoK).
\newblock In \emph{Principles of Security and Trust (POST)}, LNCS 10204, Springer, 2017.

\bibitem{BartolettiZunino2018}
M. Bartoletti, R. Zunino.
\newblock BitML: a calculus for Bitcoin smart contracts.
\newblock In \emph{ACM Conference on Computer and Communications Security (CCS)}, 2018.

\bibitem{Daian2020}
P. Daian, S. Goldfeder, T. Kell, Y. Li, X. Zhao, I. Bentov, L. Breidenbach, A. Juels.
\newblock Flash Boys 2.0: frontrunning, transaction reordering, and consensus instability in
decentralized exchanges.
\newblock In \emph{IEEE Symposium on Security and Privacy (S\&P)}, 2020.

\bibitem{LamelaThompson2018}
P. Lamela Seijas, S. Thompson.
\newblock Marlowe: financial contracts on blockchain.
\newblock In \emph{Leveraging Applications of Formal Methods, Verification and Validation (ISoLA)},
LNCS 11247:356--375, Springer, 2018.

\bibitem{EUTXO2020}
M. M. T. Chakravarty, J. Chapman, K. MacKenzie, O. Melkonian, M. Peyton Jones, P. Wadler.
\newblock The Extended UTXO model.
\newblock In \emph{Financial Cryptography and Data Security Workshops (WTSC)}, LNCS 12063:525--539,
Springer, 2020.

\bibitem{F1R3Fi}
F1R3FLY.io.
\newblock F1R3Fi: development of DeFi using Peyton Jones--Eber--Seward contracts for financial
engineering.
\newblock Open-source software repository, F1R3FLY.io, 2026.
\newblock \url{https://github.com/F1R3FLY-io/F1R3Fi}.

\bibitem{DoddFrank2010}
United States Congress.
\newblock Dodd--Frank Wall Street Reform and Consumer Protection Act.
\newblock Pub. L. No. 111--203, 2010. (Title~VII: over-the-counter derivatives reform.)

\bibitem{EMIR2012}
European Parliament and Council.
\newblock Regulation (EU) No 648/2012 on OTC derivatives, central counterparties and trade
repositories (EMIR), 2012.

\bibitem{MiFIDII2014}
European Parliament and Council.
\newblock Directive 2014/65/EU (MiFID~II) and Regulation (EU) No 600/2014 (MiFIR), 2014.

\bibitem{BaselIII2011}
Basel Committee on Banking Supervision.
\newblock Basel III: a global regulatory framework for more resilient banks and banking systems.
\newblock Bank for International Settlements, 2010 (rev. 2011).

\bibitem{MiCA2023}
European Parliament and Council.
\newblock Regulation (EU) 2023/1114 on markets in crypto-assets (MiCA), 2023.

\end{thebibliography}

\end{document}
